% Fuente: Auxiliar 8, «Métodos de programación entera».
Considere el siguiente problema de programación entera:
\[
\begin{aligned}
\max\quad & x_{1} \;+\; 2\,x_{2} \\[0.5em]
\text{sujeto a}\quad 
& -\,3\,x_{1} \;+\; 4\,x_{2} \;\le 4\\
& 3\,x_{1} \;+\; 2\,x_{2} \;\le 11\\
& 2\,x_{1} \;-\; x_{2} \;\le 5\\
& x_{1},\,x_{2} \;\ge 0,\quad x_{1},\,x_{2} \in \mathbb{Z}.
\end{aligned}
\]
Usar un gráfico para responder las siguientes preguntas:

\begin{enumerate}
  \item[(a)] ¿Cuál es el costo óptimo de la relajación lineal? ¿Y cuál es el costo óptimo del problema de programación entera?
  \item[(b)] ¿Cuál es la envoltura convexa del conjunto de todas las soluciones enteras factibles?
  \item[(c)] Ilustrar cómo funciona el algoritmo de planos de corte de Gomory. Realizar el primer corte.
  \item[(d)] Resolver el problema mediante \emph{branch \& bound}, resolviendo gráficamente las relajaciones.
  \item[(e)] Suponga que se dualiza la restricción \( -\,3\,x_{1} + 4\,x_{2} \le 4 \). ¿Cuál es el valor óptimo \(Z_{D}\) del dual lagrangeano? ¿Cómo cambia el valor de $Z_D$ si se dualiza la restricción \( 2\,x_{1} - x_{2} \le 5 \)? 
\end{enumerate}
